\chapter{Alimentation}

\section{Introduction}
La maquette fut alimenté jusqu'à maintenant pas une alimentation de laboratoire. Afin de rendre la maquette autonome, un convertisseur AC/DC fut implémenter.

\section{Vérification de l'alimentation}
L'alimentation fut dans un premier temps vérifier afin de garantir la bonne tension d'alimentation. Un potentiomètre permet de correctement définir la tension de sortie de celui-ci. Il fut dès alors fixer à 48V :



Le courant de sortie max de l'alimentation est de 8.3A.

\section{Ajout d'un fusible}
La maquette ne possède aucun fusible. Il fut dès lors ajouter afin de protéger le restant de la maquette dans le cas d'un courant bien trop éléver :


Le fusible fut fixer a 5A pour garantir un courant suffisant pour l'alimentation intégrale de la maquette et laisser une certaine marge de fonctionnemnt.

\section{Conclusion}
La maquette fut dès lors essayer avec l'alimentation. Celle-ci fonctionne sans problème venant dès lors confirmer le bon fonctionnement de l'alimentation.